\documentclass[11pt]{article}

% --- compact global formatting (no content changes) ---
\usepackage[margin=1in]{geometry}
\usepackage{microtype}
\usepackage{mathtools,amssymb}
\usepackage{tikz-cd}
\usepackage[most]{tcolorbox}
\tcbuselibrary{theorems}

% tighter display-math spacing
\setlength{\abovedisplayskip}{0.6em}
\setlength{\belowdisplayskip}{0.6em}
\setlength{\abovedisplayshortskip}{0.4em}
\setlength{\belowdisplayshortskip}{0.4em}

% slightly tighter line spacing; keep readable
\linespread{0.98}
\raggedbottom


 % compact theorem, proposition, corrolary and definition boxes
\newtcbtheorem[number within=section]{theorem}{Theorem}%
{enhanced, breakable,
 colback=blue!5!white, colframe=blue!75!black, fonttitle=\bfseries,
 boxsep=0.4ex, left=0.6ex, right=0.6ex, top=0.6ex, bottom=0.6ex,
 before skip=0.6em, after skip=0.6em}{th}
\newtcbtheorem[use counter from=theorem]{proposition}{Proposition}%
{enhanced, breakable,
 colback=blue!5!white, colframe=blue!75!black, fonttitle=\bfseries,
 boxsep=0.4ex, left=0.6ex, right=0.6ex, top=0.6ex, bottom=0.6ex,
 before skip=0.6em, after skip=0.6em}{prop}
\newtcbtheorem[use counter from=theorem]{lemma}{Lemma}%
{enhanced, breakable,
 colback=blue!5!white, colframe=blue!75!black, fonttitle=\bfseries,
 boxsep=0.4ex, left=0.6ex, right=0.6ex, top=0.6ex, bottom=0.6ex,
 before skip=0.6em, after skip=0.6em}{prop}
\newtcbtheorem[use counter from=theorem]{corollary}{Corollary}%
{enhanced, breakable,
 colback=blue!5!white, colframe=blue!75!black, fonttitle=\bfseries,
 boxsep=0.4ex, left=0.6ex, right=0.6ex, top=0.6ex, bottom=0.6ex,
 before skip=0.6em, after skip=0.6em}{cor}

\newtcbtheorem[use counter from=theorem]{definition}{Definition}%
{enhanced, breakable,
 colback=blue!5!white, colframe=blue!75!black, fonttitle=\bfseries,
 boxsep=0.4ex, left=0.6ex, right=0.6ex, top=0.6ex, bottom=0.6ex,
 before skip=0.6em, after skip=0.6em}{def}
\newtcolorbox{mybox}[1]{% #1 = title
  enhanced, breakable,
  colback=blue!5!white, colframe=blue!75!black, fonttitle=\bfseries,
  title={#1},
  boxsep=0.4ex, left=0.6ex, right=0.6ex, top=0.6ex, bottom=0.6ex,
  before skip=0.6em, after skip=0.6em
}
\usepackage{slashed}

\usepackage{float}  
\begin{document}

\section{Basic Ring Theory}

While being the purest commutative ring(free), the polynomial ring is also important for geometry. Let's say we center $S^1$ at a cartesian frame. That means we can describe it $\textbf{algebraically}$ by 
$$
x^2+y^2-1=0
$$
This expression, of course, is not unique, simply because 
$$
2x^2+2y^2-2=0
$$
and
$$
(x^2+y^2-1)^2=0
$$
also describes it. But since all of the expressions describing the circle is considered important, we need to study \textbf{a set of polynomials}, which will be our central subject matter for this section.\\

The constituting elements  of polynomials are monomials of the form:
$$
x^\alpha=x_1^{\alpha_1}...x_n^{\alpha_n}
$$
where $\alpha=(\alpha_1,...,\alpha_n)\in \mathbb{N}^n$, which determines the monomial uniquely. Consequently, a polynomial is of the form:
$$
f=\sum_{\alpha\in \mathbb{N}^n}a_\alpha x^\alpha 
$$
where $a_\alpha=0$ for all but finitely many $\alpha$. A polynomial is uniquely determined by its coefficient $a_\alpha$. \\

Next, the usual addition and cauchy product of polynomials turn the set of polynomials into a ring. 
\begin{mybox}{warning}
Note that a polynomial is different from a polynomial function. Consider the finite ring $\mathbb{Z}/2\mathbb{Z}$    and the polynomials 
$$
x \ \ and \ \ x^2
$$
They are different as polynomials but the same as functions. 
\end{mybox}
A useful identity for polynomial rings is that 
$$
R[x_1,...,x_n]\cong R[y_1,...,y_m][z_{1},...,z_k]
$$
for any partitions of $\{x_i\} $into two group $\{y_j\}$  and  $\{z_k\}$. 
\\

By partitioning into $\{x_1,...,x_{n-1}\}$ and $\{x_n\} $, a proof by induction quite easily shows that 
$$
\text{R is an integral domain} \Rightarrow \text{$R[x_1,...,x_n]$ is an integral domain }
$$
This then quickly implies that if $R$ is an integral domain, then 
$$
deg(fg)=deg(f)+deg(g)
$$
We know that "divide and conquer" is a good strategy. In order to apply this technique, we need to investigate the "atoms" of polynomials, namely their irreducible parts. \\

\textbf{Irreducibility} should mean  \textbf{"can't be broken into smaller non-trivial pieces".}  But what should be considered trivial and vice versa? We observe that units, namely elements with multiplicative inverses, don't really make something "bigger". 
\begin{mybox}{An example}
Consider the ring $\mathbb{Z}$ and the unit $1$, we see that 
$$
6=2*3=1*2*3
$$
Well, the unit here doesn't really tell us anything new about the structure of $6$. More, generally, let $a$ be a unit s.t. 
$$
p=ab
$$
$$
\Rightarrow a^{-1}p=b
$$
which says $b$ is merely a relabeling of $p$.  A genuine piece should change the factorization structure in a \textbf{non-reversible} way when we remove it.
\end{mybox}

This example implies that units are trivial, and correspondingly, non-units are non-trivial. Hence, an element $p$ being irreducible  should immediately mean that itself is not a unit nor $0$, and 
$$
p=ab\Rightarrow \text{a or b is a unit}
$$
Namely,  $p$ is irreducible if it is non trivial and cannot be written as a product of two non-trivial parts.\\

Let's compute the unit of polynomial ring. 
\begin{mybox}{Units of polynomial ring over a field}
Let $K$ be a field, then the units of $K[x_1,...,x_n]$ is the set $K/\{0\}$. To see this, take a unit $f$, then inverse exists s.t.
$$
f*f^{-1}=1
$$
Since $K$ is an integral domain, we apply the degree function on both sides
$$
deg(f)+deg(f^{-1})=0
$$
But since degree is a non-negative function, we must have 
$$
deg(f)=0
$$
This means $f$ is a constant polynomial.  Also because $f$ is invertible, $f\neq 0$.  
\end{mybox}
This example implies that irreducibility in $K[x_1,...,x_n]$ means \textbf{being non-constant polynomial, and cannot be writte as a product of two non-constant polynomials.}\\

Next, let's see what an irreducible might look like in polynomial ring

\begin{mybox}{Linear polynomial}
A polynomial $f$ with $deg(f)=1$ is necessarily irreducible because if 
$$
f=gh
$$
we have 
$$
deg(f)=deg(g)+deg(h)=1
$$
which means at least one of $g$ or $h$ is a unit. 
\end{mybox}
Then, applying a strong induction, we can prove that 
$$
K[x_1,...,x_n] \text{ is a factoriazation domain}
$$
We can actually add \textbf{uniqueness } into this characterization(which is our goal for this section). For this purpose, we investigate ( for integral domain) a special kind of atoms , called the primes. \\

A prime is a non-trivial piece that cannot hide inside a product. It cannot appear only from the "interaction" of two pieces. We see that \textbf{for a FD} , 
$$
\text{irreducibility = prime}\Rightarrow UFD
$$
Hence,  our goal is to prove that $\text{irreducibility = prime}$ in $K[x_1,...,x_n]$ . But first, we need to \textbf{develop some structural techniques in compensating for the elementwise techniques}. \\

For this purpose,  we recall that an ideal $I$ is a stable substructure of a ring $R$. Thanks to this stability, imposing the relation $I=0$ upon $R$ through the quotient $R/I$  won't disturb the existing ring structure, which means $R/I$ is a ring itself. \\

In addition to the first isomorphism theorem, there are several important structural techniques unique to rings:
\begin{mybox}{Three techniques}
\begin{enumerate}
    \item $p$ is a prime $\Leftrightarrow $ $<p>$ is a prime ideal  (This is a link between elementwise and structural techniques) 
    \item $I$ is a prime ideal $\Leftrightarrow$ $R/I$ is an integral domain 
    \item $I$ is a maximal ideal $\Leftrightarrow$ $R/I$ is a field 
\end{enumerate}    
\end{mybox}
We apply these to some calculations:

\begin{mybox}
Prove $<x_1-a_1,...,x_n-a_n>$ is a maximal ideal of $K[x_1,...,x_n]$. \\

Consider the evaluation map 

$$
ev_a: f(x_1,...,x_n)\rightarrow f(a_1,...,a_b)
$$
which is quite obviously surjective. To compute the kernel, note that for each generator $x_i-a_i$ we have 

$$
ev(x_i-a_i)=0
$$
which implies 
$$
I\subseteq ker(ev_a)
$$
The other direction can be solved using a telescoping decomposition 
$$
f(x_1,...,x_n)-f(a_1,...,a_n)=\sum_{i=1}^n(f(a_1,...,a_{i-1},x_i,...,x_n)-f(a_1,...,a_{i},x_{i+1},...,x_n))
$$
Define the one-variable polynomial
\[
G_i(t)\;:=\; f(a_1,\dots,a_{i-1},\,t,\,x_{i+1},\dots,x_n)\ \in\ K[x_{i+1},\dots,x_n][t].
\]
Then
\[
f(a_1,\dots,a_{i-1},\,x_i,\dots,x_n)\;-\;f(a_1,\dots,a_i,\,x_{i+1},\dots,x_n)
\;=\; G_i(x_i)\;-\;G_i(a_i).
\]
Now apply the factor theorem in the variable \(t\) to get
\[
G_i(x_i)\;-\;G_i(a_i)\;=\;(x_i-a_i)\,k_i(x),
\]
for some \[(k_i(x)\in K[x_1,\dots,x_n]. \] Denote

$$
g(x)=(x_i-a_i)k_i(x)
$$
We have 
$$
f(x_1,...,x_n)-f(a_1,...,a_b)=\sum_{i=1}^n(x-a_i)k_i(x)
$$
QED. Then, by the first isomorphism theorem
    $$
    \frac{K[x_1,...,x_n]}{<x_1-a_1,...,x_n-a_n>}\cong K
    $$

But because $K$ is a field, by the third technique, we get the result. 


\end{mybox}


\begin{mybox}{Three techniques }
We prove that $<x^2+1>$ is a maximal in $\mathbb{R}[x]$ by noting that 
$$
\frac{\mathbb{R}[x]}{<x^2+1>}\cong \mathbb{C}
$$
Induced by the surjective map 
$$
f(x)\rightarrow f(i)
$$

Similarly, we can prove  $<x-1>$ is a prime ideal in $\mathbb{Z}[x]$ by noting that 

$$
\frac{\mathbb{Z}[x]}{<x-1>}\cong \mathbb{Z}
$$


\end{mybox}
Now we apply them to actually prove something non-trivial! 

\begin{mybox}{A non-trivial application}
let's apply the techniques to prove the uniqueness part of 
$$
PID\Rightarrow UFD
$$
Well, this is equivalent to prove irreducibility $=$ prime in $PID$. Let's take an irreducible element $p$ and an ideal $I$ s.t. 
$$
(p)\subseteq I\subseteq R
$$
Since $R$ is a $PID$, there exists some $a\in R$ s.t.  $I=(a)$. This implies that 
$$
a|p
$$
Equivalently,  $p=ab$ for some $b\in R$.  Since $p$ is irreducible,

\begin{enumerate}
    \item If $a$ is a unit, $I=R$
    \item If $b$ is a unit, $I=(p)$
\end{enumerate}

In either case, we have $(p)$ is a maximal ideal. Hence, it is also a prime ideal, which means $p$ is a prime. \\
\end{mybox}

Next, we can apply this proposition to prove $K[x]$ is a $UFD$ by simply showing that it is a $PID$. To prove it is a $PID$, we simply apply the standard minimality argument on Euclidean function(\textbf{need more investigation here)}. \\


\end{document}
